\chapter{Complexity — Diagonal}
%%% generated by the blueprint pipeline from TCSlib.Complexity.TimeHierarchy.Diagonal
\section{Overview}

Declarations of the Lean module \texttt{TCSlib.Complexity.TimeHierarchy.Diagonal}.

\section{Declarations}

\begin{definition}[Diagonal language time bound]
\lean{Complexity.TimeHierarchy.diagLang_mem_DTIME}
\label{Complexity.TimeHierarchy.diagLang_mem_DTIME}
\leanok
\uses{Complexity.DTIME, Complexity.TimeConstructible, Complexity.TimeHierarchy.diagLang, Complexity.TimeHierarchy.diagLangOf_mem_DTIME, Complexity.TimeHierarchy.diagSim}
\difficulty{1}
Let $g:\mathbb{N}\to\mathbb{N}$ be time-constructible. The language of binary strings
$x$ for which the diagonal simulator does not halt with output $1$ after exactly
$g(|x|)$ steps is decidable by a Turing machine within time $c\bigl(g(n)+1\bigr)$ on
inputs of length $n$, for some constant $c\in\mathbb{N}$.
\end{definition}

\begin{definition}[Quadratic overhead bound]
\lean{Complexity.TimeHierarchy.diag_cost_le}
\label{Complexity.TimeHierarchy.diag_cost_le}
\leanok
\difficulty{4}
For all natural numbers $n,T,c_0,c_1,C$, we have
\[
5n+7+C\bigl(c_1(c_0T+1)^2+1\bigr) \le \bigl(8+C(c_1(c_0+1)^2+1)\bigr)(T+n+1)^2.
\]
\end{definition}

\begin{theorem}[Diagonal language outside a time class]
\lean{Complexity.TimeHierarchy.diagLang_not_mem_DTIME}
\label{Complexity.TimeHierarchy.diagLang_not_mem_DTIME}
\leanok
\uses{Complexity.DTIME, Complexity.TimeHierarchy.DiagSimulator.diagLang_not_mem_DTIME, Complexity.TimeHierarchy.diagLang, Complexity.TimeHierarchy.sqDiagSimulator}
\difficulty{1}
Let $g,T:\mathbb{N}\to\mathbb{N}$ satisfy the condition that for every
$A,N\in\mathbb{N}$ there is an $n\ge N$ such that $A\bigl(T(n)+n+1\bigr)^2\le g(n)$.
Then the language of binary strings $x$ for which the diagonal simulator does not halt
with output $1$ after exactly $g(|x|)$ steps does not belong to $\mathrm{DTIME}(T)$.
\end{theorem}

\begin{theorem}[Deterministic time hierarchy]
\lean{Complexity.time_hierarchy}
\label{Complexity.time_hierarchy}
\leanok
\uses{Complexity.DTIME, Complexity.TimeConstructible, Complexity.TimeHierarchy.DiagSimulator.time_hierarchy, Complexity.TimeHierarchy.le_sq_overhead, Complexity.TimeHierarchy.sqDiagSimulator}
\difficulty{1}
Let $f,g:\mathbb{N}\to\mathbb{N}$, and suppose $g$ is time constructible. Suppose also
that for every $A\in\mathbb{N}$ there is an $N\in\mathbb{N}$ such that for every $n\ge
N$, $A(f(n)+n+1)^2\le g(n)$. Then $\mathrm{DTIME}(f)$ is a proper subset of
$\mathrm{DTIME}(n\mapsto g(n)+1)$.
\end{theorem}

\begin{theorem}[Time hierarchy for positive bounds]
\lean{Complexity.time_hierarchy_of_pos}
\label{Complexity.time_hierarchy_of_pos}
\leanok
\uses{Complexity.DTIME, Complexity.TimeConstructible, Complexity.TimeHierarchy.DiagSimulator.time_hierarchy_of_pos, Complexity.TimeHierarchy.le_sq_overhead, Complexity.TimeHierarchy.sqDiagSimulator}
\difficulty{1}
Let $f,g:\mathbb{N}\to\mathbb{N}$, and suppose $g$ is time constructible and $g(n)>0$
for every $n$. Suppose further that for every $A\in\mathbb{N}$ there is an
$N\in\mathbb{N}$ such that for every $n\ge N$, $A(f(n)+n+1)^2\le g(n)$. Then
$\mathrm{DTIME}(f)$ is a proper subset of $\mathrm{DTIME}(g)$.
\end{theorem}

%%% added by the blueprint pipeline
\section{Additional declarations}

\begin{definition}[Fixed effective machine coding]
\lean{Complexity.TimeHierarchy.code}
\label{Complexity.TimeHierarchy.code}
\leanok
\uses{Turing.EffectiveMachineCode, Turing.exists_effectiveMachineCode}
Fix an effective machine coding: an encoding of coded Boolean Turing machines by Boolean
words and a total decoding map such that every machine is recovered when its code is
followed by any number of padding bits equal to $1$. The coding also includes a finite
Boolean machine that computes the serialization of the machine decoded from any input
word, together with a time bound for that computation.
\end{definition}

\begin{definition}[Universal Boolean Turing machine]
\lean{Complexity.TimeHierarchy.univTM}
\label{Complexity.TimeHierarchy.univTM}
\leanok
\uses{Complexity.TimeHierarchy.code, Turing.FinTM, Turing.universal}
For the fixed effective machine coding, choose a finite Boolean Turing machine $U$ that
simulates every decoded machine with the following guarantees. For each code word
$\alpha$, there is a constant $C$ such that every computation of the decoded machine on
input $x$ that produces output $y$ in time $t$ is simulated by $U$ on the paired
encoding of $\alpha$ and $x$, producing $y$ in time $C(t+1)$. Conversely, whenever $U$
on that paired input produces an output, the decoded machine on $x$ produces the same
output in some finite time.
\end{definition}

\begin{lemma}[Universal simulation of coded machines]
\lean{Complexity.TimeHierarchy.univTM_spec}
\label{Complexity.TimeHierarchy.univTM_spec}
\leanok
\uses{Complexity.TimeHierarchy.code, Complexity.TimeHierarchy.univTM, Turing.CodeTM.toFinTM, Turing.FinTM.ComputesInTime, Turing.pairEncode, Turing.universal}
\difficulty{1}
For every bit string $\alpha$, there is a constant $C\in\mathbb{N}$ such that for every
bit string $x$, both of the following hold: if the machine decoded from $\alpha$
computes an output $y$ from $x$ within $t$ steps, then the universal machine computes
$y$ from the paired encoding of $\alpha$ and $x$ within $C(t+1)$ steps; and if the
universal machine computes some output $y$ from that paired encoding within some finite
number of steps, then the machine decoded from $\alpha$ computes $y$ from $x$ within
some finite number of steps.
\end{lemma}

\begin{definition}[Diagonal simulation machine]
\lean{Complexity.TimeHierarchy.diagSim}
\label{Complexity.TimeHierarchy.diagSim}
\leanok
\uses{Complexity.TimeHierarchy.preTM, Complexity.TimeHierarchy.univTM, Turing.FinTM, Turing.FinTM.bufferedCompTM}
A particular finite Turing machine over the Boolean alphabet, obtained by buffering the
pre-universal machine with the universal machine.
\end{definition}

\begin{definition}[Diagonal language]
\lean{Complexity.TimeHierarchy.diagLang}
\label{Complexity.TimeHierarchy.diagLang}
\leanok
\uses{Complexity.TimeHierarchy.diagSim, Turing.FinTM.ComputesInTime}
For a function $g : \mathbb{N} \to \mathbb{N}$, define a language of finite binary
strings by declaring $x$ to belong to it exactly when the diagonal simulation machine
does not compute the one-bit output $1$ on input $x$ within $g(|x|)$ steps.
\end{definition}

\begin{definition}[Diagonal language for a machine]
\lean{Complexity.TimeHierarchy.diagLangOf}
\label{Complexity.TimeHierarchy.diagLangOf}
\leanok
\uses{Turing.FinTM, Turing.FinTM.ComputesInTime}
Given a finite Turing machine $W$ over the Boolean alphabet and a function $g :
\mathbb{N} \to \mathbb{N}$, define a language of finite binary strings by declaring $x$
to belong to it exactly when $W$ does not compute the one-bit output $1$ on input $x$
within $g(|x|)$ steps.
\end{definition}

\begin{theorem}[Equality of diagonal languages]
\lean{Complexity.TimeHierarchy.diagLang_eq_diagLangOf}
\label{Complexity.TimeHierarchy.diagLang_eq_diagLangOf}
\leanok
\uses{Complexity.TimeHierarchy.diagLang, Complexity.TimeHierarchy.diagLangOf, Complexity.TimeHierarchy.diagSim}
\difficulty{1}
For every function $g:\mathbb{N}\to\mathbb{N}$, the language of bit strings $x$ for
which the diagonal simulation machine does not halt within the specified time $g(|x|)$
with output equal to the one-bit string $1$ is equal to the language of bit strings $x$
for which the diagonal simulation does not halt within that time with that output.
\end{theorem}

\begin{theorem}[Diagonal language time bound]
\lean{Complexity.TimeHierarchy.diagLangOf_mem_DTIME}
\label{Complexity.TimeHierarchy.diagLangOf_mem_DTIME}
\leanok
\uses{Complexity.DTIME, Complexity.TimeConstructible, Complexity.TimeHierarchy.clockTM, Complexity.TimeHierarchy.clockTM_spec, Complexity.TimeHierarchy.ctrVal_bits, Complexity.TimeHierarchy.diagLangOf, Turing.FinTM, Turing.FinTM.ComputesInTime, Turing.FinTM.ComputesInTime.mono, Turing.MultiTapeTM.indicator, Turing.length_bits_le_self}
\difficulty{5}
Let $W$ be a finite-state Turing machine with Boolean input and output, and let
$g:\mathbb{N}\to\mathbb{N}$ be time-constructible. The language consisting of words $x$
for which $W$ does not halt at time $g(|x|)$ with output the single bit $1$ is decidable
in time $c\bigl(g(n)+1\bigr)$ for some constant $c\in\mathbb{N}$.
\end{theorem}

\begin{definition}[Padded diagonal simulator]
\lean{Complexity.TimeHierarchy.DiagSimulator}
\label{Complexity.TimeHierarchy.DiagSimulator}
\leanok
\uses{Complexity.DTIME, Turing.FinTM, Turing.FinTM.ComputesInTime, Turing.MultiTapeTM.indicator}
For a function $\Phi : \mathbb{N} \times \mathbb{N} \to \mathbb{N}$, a padded diagonal
simulator consists of a finite Turing machine $W$ over the Boolean alphabet with the
following property: for every time bound $T : \mathbb{N} \to \mathbb{N}$ and every
language $L \in \mathrm{DTIME}(T)$, there are natural numbers $A$ and $N_0$ such that,
for every $n \ge N_0$, some binary string $x$ of length $n$ is computed by $W$ with
output the indicator of whether $x \in L$, within $A\Phi(T(n),n)$ steps.
\end{definition}

\begin{definition}[Monotonicity of diagonal simulation bounds]
\lean{Complexity.TimeHierarchy.DiagSimulator.mono}
\label{Complexity.TimeHierarchy.DiagSimulator.mono}
\leanok
\uses{Complexity.TimeHierarchy.DiagSimulator, Turing.FinTM.ComputesInTime.mono}
Suppose a padded diagonal simulator works with bound function $\Phi$, and let $\Phi' :
\mathbb{N} \times \mathbb{N} \to \mathbb{N}$ and $K \in \mathbb{N}$ satisfy $\Phi(t,n)
\le K\Phi'(t,n)$ for all natural numbers $t,n$. Then the same simulating machine is a
padded diagonal simulator with bound function $\Phi'$.
\end{definition}

\begin{theorem}[Diagonal language escapes the time class]
\lean{Complexity.TimeHierarchy.DiagSimulator.diagLang_not_mem_DTIME}
\label{Complexity.TimeHierarchy.DiagSimulator.diagLang_not_mem_DTIME}
\leanok
\uses{Complexity.DTIME, Complexity.TimeHierarchy.DiagSimulator, Complexity.TimeHierarchy.diagLangOf, Turing.FinTM.ComputesInTime, Turing.FinTM.ComputesInTime.mono, Turing.FinTM.ComputesInTime.output_unique, Turing.MultiTapeTM.indicator}
\difficulty{5}
Let $S$ be a diagonal simulator for a function
$\Phi:\mathbb{N}\times\mathbb{N}\to\mathbb{N}$, and let $g,T:\mathbb{N}\to\mathbb{N}$.
Suppose that for every natural number $A$ and every threshold $N$, there is some $n\ge
N$ such that $A\Phi(T(n),n)\le g(n)$. Then the language consisting of inputs on which
the simulator's machine does not compute the one-symbol output $1$ within $g(n)$ steps,
where $n$ is the input length, does not belong to $\mathrm{DTIME}(T)$.
\end{theorem}

\begin{theorem}[Deterministic time hierarchy]
\lean{Complexity.TimeHierarchy.DiagSimulator.time_hierarchy}
\label{Complexity.TimeHierarchy.DiagSimulator.time_hierarchy}
\leanok
\uses{Complexity.DTIME, Complexity.TimeConstructible, Complexity.TimeHierarchy.DiagSimulator, Complexity.TimeHierarchy.DiagSimulator.diagLang_not_mem_DTIME, Complexity.TimeHierarchy.diagLangOf_mem_DTIME, Turing.FinTM.ComputesInTime.mono}
\difficulty{5}
Let $S$ be a diagonal simulator for $\Phi:\mathbb{N}\times\mathbb{N}\to\mathbb{N}$.
Assume $t\le\Phi(t,n)$ for all natural numbers $t,n$, and let
$f,g:\mathbb{N}\to\mathbb{N}$ with $g$ time-constructible. Suppose that for every
natural number $A$ there is a threshold $N$ such that $A\Phi(f(n),n)\le g(n)$ for all
$n\ge N$. Then $\mathrm{DTIME}(f)$ is a proper subset of $\mathrm{DTIME}(g+1)$.
\end{theorem}

\begin{theorem}[Time hierarchy for positive bounds]
\lean{Complexity.TimeHierarchy.DiagSimulator.time_hierarchy_of_pos}
\label{Complexity.TimeHierarchy.DiagSimulator.time_hierarchy_of_pos}
\leanok
\uses{Complexity.DTIME, Complexity.DTIME.mono, Complexity.TimeConstructible, Complexity.TimeHierarchy.DiagSimulator, Complexity.TimeHierarchy.DiagSimulator.time_hierarchy, Turing.FinTM.ComputesInTime.mono}
\difficulty{4}
Let $S$ be a diagonal simulator for $\Phi:\mathbb{N}\times\mathbb{N}\to\mathbb{N}$.
Assume $t\le\Phi(t,n)$ for all natural numbers $t,n$, and let
$f,g:\mathbb{N}\to\mathbb{N}$ with $g$ time-constructible and $g(n)>0$ for every $n$.
Suppose that for every natural number $A$ there is a threshold $N$ such that
$A\Phi(f(n),n)\le g(n)$ for all $n\ge N$. Then $\mathrm{DTIME}(f)$ is a proper subset of
$\mathrm{DTIME}(g)$.
\end{theorem}

\begin{definition}[Self-simulation structure]
\lean{Complexity.TimeHierarchy.SelfSimulator}
\label{Complexity.TimeHierarchy.SelfSimulator}
\leanok
\uses{Complexity.DTIME, Turing.FinTM, Turing.FinTM.ComputesInTime, Turing.MultiTapeTM.indicator, Turing.pairEncode}
For a function $\Psi : \mathbb{N} \to \mathbb{N}$, a self-simulation structure consists
of a finite Boolean Turing machine $W$ and a decoder assigning a finite Boolean Turing
machine to each binary string $\alpha$. For each code $\alpha$, there is a constant $C$
such that whenever the machine decoded from $\alpha$ computes a one-bit output $b$ on
the paired input $\langle\alpha,w\rangle$ within $t$ steps, $W$ computes the same output
on that input within $C(t+|\langle\alpha,w\rangle|+1)$ steps. Moreover, for every time
bound $T$ and every $L \in \mathrm{DTIME}(T)$, there are a code $\alpha$ and a constant
$c$ such that, on every input $x$, the machine decoded from $\alpha$ computes the
indicator of membership in $L$ within $c(\Psi(T(|x|))+1)$ steps.
\end{definition}

\begin{definition}[Diagonal simulator from self-simulation]
\lean{Complexity.TimeHierarchy.SelfSimulator.toDiagSimulator}
\label{Complexity.TimeHierarchy.SelfSimulator.toDiagSimulator}
\leanok
\uses{Complexity.TimeHierarchy.DiagSimulator, Complexity.TimeHierarchy.SelfSimulator, Turing.FinTM.ComputesInTime.mono, Turing.length_pairEncode, Turing.pairEncode}
Given a self-simulation structure with bound function $\Psi$, its machine $W$ is a
padded diagonal simulator with bound function $(t,n) \mapsto \Psi(t)+n+1$.
\end{definition}

\begin{definition}[One-tape self-simulation structure]
\lean{Complexity.TimeHierarchy.oneTapeSelfSimulator}
\label{Complexity.TimeHierarchy.oneTapeSelfSimulator}
\leanok
\uses{Complexity.TimeHierarchy.SelfSimulator, Complexity.TimeHierarchy.code, Complexity.TimeHierarchy.diagSim, Complexity.TimeHierarchy.preTM, Complexity.TimeHierarchy.preTM_computes, Complexity.TimeHierarchy.scanPre_pairEncode_append, Complexity.TimeHierarchy.univTM, Complexity.TimeHierarchy.univTM_spec, Turing.CodeTM.toFinTM, Turing.FinTM.ComputesInTime.mono, Turing.FinTM.bufferedCompTM_computesInTime, Turing.FinTM.one_work_tape_binary, Turing.MachineCode.decode_encode, Turing.MultiTapeTM.indicator, Turing.exists_codeTM, Turing.length_pairEncode, Turing.pairEncode}
A particular self-simulation structure is defined with overhead function $t \mapsto
(t+1)^2$. Its simulating machine is the diagonal simulation machine, and its decoder
maps a binary code to the corresponding coded Turing machine. For each code, it
simulates that machine's one-bit computations on paired inputs with a constant-factor
overhead in the computation time plus input length; it also represents every language in
$\mathrm{DTIME}(T)$ by a code whose decoded machine decides the language on every input
within $c(T(|x|)+1)^2$ steps for some constant $c$.
\end{definition}

\begin{theorem}[Quadratic overhead dominates time]
\lean{Complexity.TimeHierarchy.le_sq_overhead}
\label{Complexity.TimeHierarchy.le_sq_overhead}
\leanok
\difficulty{3}
For all natural numbers $t$ and $n$, $t \le (t+n+1)^2$.
\end{theorem}

\begin{definition}[Quadratic-Time Diagonal Simulator]
\lean{Complexity.TimeHierarchy.sqDiagSimulator}
\label{Complexity.TimeHierarchy.sqDiagSimulator}
\leanok
\uses{Complexity.TimeHierarchy.DiagSimulator, Complexity.TimeHierarchy.DiagSimulator.mono, Complexity.TimeHierarchy.SelfSimulator.toDiagSimulator, Complexity.TimeHierarchy.oneTapeSelfSimulator}
There exists a finite Turing machine over the binary alphabet that simulates every
language in $\mathrm{DTIME}(T)$ on inputs of length $n$ within a constant multiple of
$(T(n)+n+1)^2$ steps for all sufficiently large $n$: for every time bound $T$ and every
such language $L$, if $L\in\mathrm{DTIME}(T)$, then there are natural numbers $A$ and
$N_0$ such that for every $n\ge N_0$ some binary string $x$ of length $n$ is decided by
the machine in $A(T(n)+n+1)^2$ steps.
\end{definition}

\begin{theorem}[Diagonal simulator component identity]
\lean{Complexity.TimeHierarchy.sqDiagSimulator_W}
\label{Complexity.TimeHierarchy.sqDiagSimulator_W}
\leanok
\uses{Complexity.TimeHierarchy.diagSim, Complexity.TimeHierarchy.sqDiagSimulator}
\difficulty{1}
The component $W$ of the square diagonal simulator is equal to the machine obtained by
composing the preliminary machine with the universal machine using buffered composition.
\end{theorem}
